% ============================================================================
% SECTION I — INTRODUCTION           target 6–8 pages, 2600–3600 words
%
% PURPOSE (Bible §8): establish that the CHEMICAL question is real, unsolved and
% worth a Chemistry panel's attention; demonstrate literature thoroughness (C6);
% stake the novelty claim (C2).
%
% FRAMING RULE (anti-pattern 1): this is a COORDINATION CHEMISTRY paper, not an
% environmental engineering paper. The stated object of study is
%   "why does a d10s2 post-transition ion outrank a d9 ion at an O-donor site?"
% with remediation as the application. Referees for the Chemistry panel reward
% chemical mechanism; they discount "we made an adsorbent".
% ============================================================================
\section{Introduction}
\label{sec:introduction}

% ==========================================================================
% CHECKLIST — Section I
%   [ ] Frame as coordination chemistry in the FIRST paragraph, not as remediation
%   [ ] Cite Irving & Williams (1953) explicitly by name in §1.2
%   [ ] Four literature strands, each ENDING in an explicit gap sentence (§1.3)
%   [ ] Novelty box in §1.4 with the numbered original contributions
%   [ ] Novelty claim 1 must match reality: ossein was PURCHASED, not prepared (audit B10.6)
%   [ ] 35-60 references overall; §1.3 carries most of them
%   [ ] Every DOI verified via scripts/verify_dois.py before any citation is written
% ==========================================================================

\subsection{Trace \PbII{} in mixed-metal aqueous streams}
\label{sec:problem}
% 0.75–1 page. Two paragraphs of scale and toxicity framing, kept SHORT and cited.
% Then: why SELECTIVE RECOVERY rather than bulk removal — recovery enables
% valorisation, reduces sludge, permits reuse. This is where the concentration
% factor becomes meaningful.
% Name the real streams with real concentration ranges AND citations:
%   battery recycling leachate · e-waste · mine drainage · electroplating rinse.

\subsection{The Irving--Williams trap in non-specific biosorbents}
\label{sec:irving-williams}
% 1–1.5 pages. THE CONCEPTUAL SEED OF THE WHOLE PAPER. Bible: "write it three
% times until it is airtight."
%
%   1. State the Irving–Williams series explicitly; cite Irving & Williams 1953.
%   2. Mechanism of the trap: finite O-donor sites; higher-stability transition
%      metal complexes outcompete for them; a non-specific sorbent is
%      THERMODYNAMICALLY STEERED AWAY from Pb(II).
%   3. CRITICAL FRAMING MOVE: Pb(II) sits OUTSIDE the classical series because it
%      is not a first-row d-block ion. Selectivity for Pb(II) is therefore not a
%      matter of beating Irving–Williams on its own terms, but of engineering a
%      site whose SELECTION RULES ARE DIFFERENT — geometry, polarisability,
%      covalency and desolvation rather than charge density.


% --- registry Fig 1.1 | source: literature values (cited) | script: figures/src/fig1\_1\_irving\_williams.py ---
\begin{figure}[htbp]
  \centering
  \NEEDSDATA{Figure~\ref{fig:irving-williams} (registry 1.1)}{literature values (cited) \\ owning script: \texttt{figures/src/fig1\_1\_irving\_williams.py}}
  \caption[The Irving--Williams series with \PbII{} placed outside it]{The Irving--Williams stability sequence for first-row divalent transition-metal complexes, with \PbII{} placed outside the series. The sequence orders complex stability for a common ligand set; \PbII{}, as a $\mathrm{d}^{10}\mathrm{s}^2$ post-transition ion, is not a member of it, and its affinity at an oxygen-donor site is therefore governed by different variables.}
  \label{fig:irving-williams}
\end{figure}

% --- registry Fig 1.2 | source: schematic; geometries from this work's DFT output | script: figures/src/fig1\_2\_holo\_hemi.py ---
\begin{figure}[htbp]
  \centering
  \NEEDSDATA{Figure~\ref{fig:holo-hemi} (registry 1.2)}{schematic; geometries from this work's DFT output \\ owning script: \texttt{figures/src/fig1\_2\_holo\_hemi.py}}
  \caption[Holodirected versus hemidirected coordination geometry]{Holodirected versus hemidirected coordination geometry at a divalent metal centre. In the holodirected case the donor atoms surround the metal evenly; in the hemidirected case they are gathered into one hemisphere, leaving a void occupied by the stereochemically active $6\mathrm{s}^2$ lone pair. The void-hemisphere angle $\theta_{\mathrm{void}}$ and the normalised centroid displacement $\tilde{d}$, defined in \cref{sec:hemidirection-method}, quantify the distinction.}
  \label{fig:holo-hemi}
\end{figure}

\subsection{Existing research}
\label{sec:literature}
% 2–2.5 pages, the literature review proper. FOUR STRANDS, each ending in an
% explicit gap sentence. This subsection carries criterion C6 on its own.

\subsubsection{Polyphenol and tannin sorbents for divalent metals}
\label{sec:lit-tannin}
% Şengil & Özacar valonia tannin resin and comparable systems; typical capacities;
% typical selectivity ordering.
% GAP SENTENCE: selectivity is reported, seldom explained at the electronic-structure level.

\subsubsection{Collagen and ossein scaffolds}
\label{sec:lit-collagen}
% Fish-scale collagen isolation and thermal stability (Pati et al.); why a
% demineralised collagen scaffold is attractive — abundant carboxylate and amide,
% fibrous, waste-derived, low cost.
% GAP SENTENCE: unfunctionalised collagen is non-selective.

\subsubsection{\PbII{} coordination and the stereochemically active $6\mathrm{s}^2$ lone pair}
\label{sec:lit-hemidirection}
% Shimoni-Livny, Glusker & Bock on holodirected vs hemidirected geometries; the
% void-hemisphere concept; when hemidirection is favoured — low coordination
% number, soft/polarisable donors, anionic ligands.
% GAP SENTENCE: hemidirection is well documented in crystallography, rarely used
% as a DESIGN PRINCIPLE for a sorbent.

\subsubsection{Computational treatment of metal--polyphenol binding}
\label{sec:lit-computational}
% DFT-D3, implicit solvation (SMD), energy-decomposition and charge-decomposition
% approaches; what they have and have not been applied to.
% GAP SENTENCE: no matched experimental-competitive and decomposition study for a
% designed galloyl pocket across Pb/Cu/Zn.

\subsection{Hypothesis, aim and original contributions}
\label{sec:hypothesis}
% 1 page. Bible §10 Device 1 — the NOVELTY BOX. A referee must be able to state
% the novelty in one sentence after reading this page.

\begin{noveltybox}
\textbf{Hypothesis.}
% A vicinal-trihydroxy (galloyl) oxygen-donor pocket immobilised on a collagen
% scaffold will invert the expected Irving–Williams ordering in favour of Pb(II),
% because (i) the pocket accommodates the stereochemically active 6s2 lone pair in
% a hemidirected arrangement, (ii) donation into the diffuse Pb 6p manifold
% provides a covalent channel unavailable to the harder, more contracted Cu(II)
% and Zn(II) acceptors, and (iii) Pb(II)'s markedly lower hydration free energy
% imposes a smaller desolvation penalty on complexation.
\NEEDSDATA{Hypothesis statement}{written once §4.6 and §4.7 have produced numbers; the hypothesis as stated must be the one actually tested}

\medskip
\textbf{Aim.}
% To test this hypothesis by matched experiment and first-principles calculation,
% and to determine whether the observed selectivity is electrostatically or
% covalently/desolvation controlled.

\medskip
\textbf{Original contributions of this work.}
\begin{enumerate}\setstretch{1.0}
  \item % REWORDED per audit B10.6 / Phase 4 answer Q5: the ossein was PURCHASED
        % from Nizona Marine Products, not prepared by the author. The material is
        % waste-derived but the waste stream was processed by the supplier. State
        % it that way. Overstating provenance unravels under oral examination.
        \TODOPAL{Confirm the wording of contribution 1 now that the ossein is known to be commercially purchased. Proposed: "Synthesis and full characterisation of a galloyl-functionalised biosorbent built on a commercially sourced, waste-derived fish-scale ossein scaffold."}
  \item Quantitative demonstration of \PbII{} selectivity under a deliberately adverse molar ratio,
        with \PbII{} at a \PENDING{molar_ratio_cu_to_pb_minority}{molar deficit to Cu in the minority-target ternary}-fold molar deficit.
  \item A DFT model that reproduces the experimental ordering without parameter fitting.
  \item An energy decomposition that falsifies the purely electrostatic rationale and localises the
        selectivity to orbital covalency and desolvation.
  \item Demonstration that the mechanistic finding translates to a working, regenerable fixed-bed
        process with a \PENDING{ef_cycle1}{measured enrichment factor, cycle 1}-fold concentration factor.
\end{enumerate}
\end{noveltybox}

\subsection{Structure of this report}
\label{sec:structure}
% 0.25 page. One short paragraph mapping the sections. Costs three sentences and
% orients the referee. Winners do this.
